% English translation of chapter3.tex lines 57--162.
% Source: Methods of Algebra, Volume 2, commit
% 9a5803ff77dd3257484cb177f851a73770a59dd3 (CC BY 4.0).
% stable-unit-id: o014.aljabr2.chapter3.complexes-over-additive-categories
% Scope: complete Section 3.1.

% segment-id: o014.li2.chapter3.complexes-over-additive-categories.h001
\section{Complexes over Additive Categories}\label{sec:additive-cplx}
\hypertarget{o014.aljabr2.chapter3.complexes-over-additive-categories}{ }

% segment-id: o014.li2.chapter3.complexes-over-additive-categories.p001
Definition~\ref{def:complex} briefly introduced complexes and their usual
notation. Let $\mathcal{A}$ be an additive category, and consider unbounded
complexes over it, that is, complexes without endpoint terms,
$X=(X^n,d^n)_{n\in\Z}$. The morphism $d^n$ is often denoted by $d_X^n$.
For historical reasons, the morphisms $d_X^n$ are also called
``differentials.''
\index{differential}

% segment-id: o014.li2.chapter3.complexes-over-additive-categories.d001
\begin{definition}\label{def:cplx-cat}
	\index{complex}
	\index[sym1]{CA@$\cate{C}(\mathcal{A})$}
	A morphism from a complex $X$ to a complex $Y$ is defined to be data
	$(f^n)_{n\in\Z}$, also abbreviated to $f$, where
	$f^n\in\Hom_{\mathcal{A}}(X^n,Y^n)$ satisfies
% segment-id: o014.li2.chapter3.complexes-over-additive-categories.q001
	\[ d_Y^n f^n=f^{n+1}d_X^n,\quad n\in\Z. \]

	All complexes and their morphisms form a category $\cate{C}(\mathcal{A})$.
	The identity morphism of a complex $X$ is determined by
	$(\identity_X)^n:=\identity_{X^n}$, and composition is defined by
	$(gf)^n=g^nf^n$. Addition
	$f+g:=(f^n+g^n)_{n\in\Z}$ makes
	$\Hom_{\cate{C}(\mathcal{A})}(X,Y)$ an abelian group.
\end{definition}

% segment-id: o014.li2.chapter3.complexes-over-additive-categories.p002
If $\mathcal{A}$ is a $\Bbbk$-linear category in the sense of
Definition~\ref{def:k-linear-cat}, for a fixed commutative ring $\Bbbk$,
then $\Hom_{\cate{C}(\mathcal{A})}(X,Y)$ is also a $\Bbbk$-module under
termwise scalar multiplication. In every result of this section, ``additive''
may be replaced by the more general term ``$\Bbbk$-linear.'' For brevity,
only the additive version will be stated.

% segment-id: o014.li2.chapter3.complexes-over-additive-categories.p003
The key property in the definition of a complex is $d^{n+1}d^n=0$. It can
also be recast in the language of differential graded objects.%
\footnote{The related discussion continues in
\S~\sourcecrossref{sec:spectral-sequence}{5.2}.}

% segment-id: o014.li2.chapter3.complexes-over-additive-categories.d002
\begin{definition}[Graded objects]\label{def:graded-obj}
	\index{graded object}
	\index{graded object!$\Z^m$-graded; bigraded}
	Let $\mathcal{A}$ be an arbitrary category. Recall the construction of the
	product category $\mathcal{A}^{\Z}$:
% segment-id: o014.li2.chapter3.complexes-over-additive-categories.l001
	\begin{compactitem}
% segment-id: o014.li2.chapter3.complexes-over-additive-categories.i001
		\item its objects are of the form $X=(X^n)_{n\in\Z}$, where
		$X^n\in\Obj(\mathcal{A})$;
% segment-id: o014.li2.chapter3.complexes-over-additive-categories.i002
		\item its morphisms are of the form
		$f=(f^n:X^n\to Y^n)_{n\in\Z}$, with no further condition, and
		composition is performed termwise (or degreewise), so that
		$(gf)^n=g^nf^n$;
% segment-id: o014.li2.chapter3.complexes-over-additive-categories.i003
		\item define an automorphism $T$ of $\mathcal{A}^{\Z}$ by the following
		``shift'' functor: $(TX)^n=X^{n+1}$ and $(Tf)^n=f^{n+1}$.
	\end{compactitem}
	An object $X=(X^n)_n$ is also called a \emph{$\Z$-graded object} over
	$\mathcal{A}$, or simply a \emph{graded object}.

	More generally, objects of $\mathcal{A}^{\Z^m}$ are called
	$\Z^m$-graded objects and are written
	$(X^{n_1,\ldots,n_m})_{n_1,\ldots,n_m}$; their morphisms are written
	$(f^{n_1,\ldots,n_m})_{n_1,\ldots,n_m}$. This category carries a family
	of mutually commuting shift functors $T_1,\ldots,T_m$. The special case
	$m=2$ is also called a \emph{bigraded object}.
\end{definition}

% segment-id: o014.li2.chapter3.complexes-over-additive-categories.p004
If $\mathcal{A}$ is additive, then so is $\mathcal{A}^{\Z^m}$: addition of
morphisms, direct sums of objects, and all other operations are performed
degreewise.

% segment-id: o014.li2.chapter3.complexes-over-additive-categories.d003
\begin{definition}[Differential graded objects]
	\index{differential graded object}
	Let $\mathcal{A}$ be an additive category. A differential graded object
	over $\mathcal{A}$ is data $(X,d)$, where
	$X\in\Obj(\mathcal{A}^{\Z})$ and $d:X\to TX$ satisfies $(Td)d=0$. A
	morphism $(X,d_X)\to(Y,d_Y)$ is a morphism
	$f\in\Hom_{\mathcal{A}^{\Z}}(X,Y)$ satisfying $(Tf)d_X=d_Yf$.
\end{definition}

% segment-id: o014.li2.chapter3.complexes-over-additive-categories.p005
Writing $d$ in a differential graded object $(X,d)$ concretely as a family
of morphisms $d_X^n:X^n\to X^{n+1}$, the condition on $d$ is
$d_X^{n+1}d_X^n=0$, while the condition on a morphism is
$f^{n+1}d_X^n=d_Y^nf^n$. We have thus recovered
Definition~\ref{def:cplx-cat}. Hence the category of differential graded
objects over $\mathcal{A}$ is isomorphic to the category of complexes. We
shall pass between these two viewpoints without further comment.

% segment-id: o014.li2.chapter3.complexes-over-additive-categories.p006
The first step in studying complexes is to understand the various limits in
$\cate{C}(\mathcal{A})$. The version for the product category
$\mathcal{A}^{\Z}$ is straightforward: limits can be formed at each
$n\in\Z$, that is, termwise or degreewise in $\mathcal{A}$. We now explain
how to lift this construction to $\cate{C}(\mathcal{A})$. Write
$U:\cate{C}(\mathcal{A})\to\mathcal{A}^{\Z}$ for the forgetful functor
sending a differential graded object $(X,d)$ to its underlying graded object
$X$.

% segment-id: o014.li2.chapter3.complexes-over-additive-categories.le001
\begin{lemma}[Termwise construction of limits]\label{prop:limits-cplx}
	For any functor $\alpha:I\to\cate{C}(\mathcal{A})$, write
	$\overline{\alpha}:=U\alpha$. Suppose that
	$\varinjlim\overline{\alpha}$ exists. It admits a unique structure as an
	object of $\cate{C}(\mathcal{A})$ for which it is $\varinjlim\alpha$.
	The corresponding statement holds for $\varprojlim$.

	In the language of Definition~\ref{def:create-limit}, this says that the
	forgetful functor $U$ creates $\varinjlim$ and $\varprojlim$.
\end{lemma}
% segment-id: o014.li2.chapter3.complexes-over-additive-categories.pf001
\begin{proof}
	Suppose that $\alpha:I\to\cate{C}(\mathcal{A})$ sends
	$i\in\Obj(I)$ to the differential graded object
	$\alpha(i)=(\overline{\alpha}(i),d_i)$. If
	$\varinjlim\overline{\alpha}$ exists, the morphisms $d_i$ determine
	$d:\varinjlim\overline{\alpha}\to
	(T\varinjlim\overline{\alpha}\simeq
	\varinjlim T\overline{\alpha})$; here the automorphism $T$ automatically
	preserves $\varinjlim$. This morphism is characterized by the commutative
	diagram
% segment-id: o014.li2.chapter3.complexes-over-additive-categories.g001
	\[\begin{tikzcd}
		\varinjlim \overline{\alpha} \arrow[r, "d"] & \varinjlim T\overline{\alpha} \arrow[r, "\sim"] & T\varinjlim \overline{\alpha} \\
		\overline{\alpha}(i) \arrow[u] \arrow[r, "{d_i}"'] & T\overline{\alpha}(i) \arrow[u] \arrow[ru] &
	\end{tikzcd} \qquad (i \in \Obj(I)) . \]

	The relations $(Td_i)d_i=0$ imply $(Td)d=0$, so
	$(\varinjlim\overline{\alpha},d)$ is an object of
	$\cate{C}(\mathcal{A})$. The diagram also shows that this is the unique
	choice making all maps
	$(\overline{\alpha}(i),d_i)\to
	(\varinjlim\overline{\alpha},d)$ morphisms.

	We verify its universal property as $\varinjlim\alpha$ in
	$\cate{C}(\mathcal{A})$. Suppose a compatible family of morphisms
	$f_i:\alpha(i)\to L$ in $\cate{C}(\mathcal{A})$ is given. Forgetting the
	differentials gives $f_i:\overline{\alpha}(i)\to U(L)$, which determines a
	unique $f:\varinjlim\overline{\alpha}\to U(L)$. The key point is to show
	that $f$ is a morphism in $\cate{C}(\mathcal{A})$. Although this is to be
	expected, we write the diagram explicitly:
% segment-id: o014.li2.chapter3.complexes-over-additive-categories.g002
	\[\begin{tikzcd}
		\varinjlim \overline{\alpha} \arrow[dd, "f"'] \arrow[rr, "d"] & & \varinjlim T\overline{\alpha} \arrow[r, "\sim"] & T\varinjlim \overline{\alpha} \arrow[dd, "{Tf}"] \\
		& \overline{\alpha}(i) \arrow[lu] \arrow[ld, "{f_i}"] \arrow[r, "d_i"] & T\overline{\alpha}(i) \arrow[ru] \arrow[rd, "{Tf_i}"] \arrow[u] & \\
		U(L) \arrow[rrr, "d_L"'] & & & TU(L)
	\end{tikzcd} \qquad (i \in \Obj(I)). \]
	Commutativity of every part follows directly from the definitions or the
	construction. Thus the composites of $d_Lf$ and $(Tf)d$ with every map
	$\overline{\alpha}(i)\to\varinjlim\overline{\alpha}$ agree. Therefore
	$d_Lf=(Tf)d$, as required.
\end{proof}

% segment-id: o014.li2.chapter3.complexes-over-additive-categories.pr001
\begin{proposition}\label{prop:additive-cat-cplx}
	The category $\cate{C}(\mathcal{A})$ is additive. If $\mathcal{A}$ is
	complete (respectively, cocomplete), then $\cate{C}(\mathcal{A})$ is also
	complete (respectively, cocomplete).
\end{proposition}
% segment-id: o014.li2.chapter3.complexes-over-additive-categories.pf002
\begin{proof}
	Use Lemma~\ref{prop:limits-cplx} to reduce every required limit to a
	degreewise limit in $\mathcal{A}$.
\end{proof}

% segment-id: o014.li2.chapter3.complexes-over-additive-categories.d004
\begin{definition}[Translation functors]\label{def:translation-functor}
	\index{translation functor}
	\index[sym1]{[n]@$[n]$}
	Let $X$ be an object of $\cate{C}(\mathcal{A})$ and let $n\in\Z$. Define
	the complex $X[n]$ as follows:%
	\footnote{Inserting $(-1)^n$ in the definition of $d_{X[n]}$ is useful.
	This version is also isomorphic to the signless version
	$X[n]^\circ:=\left(X^{n+k},d_X^{k+n}\right)_k$. For $n=1$, define
	$X[1]\rightiso X[1]^\circ$ using the sign pattern
	$\cdots,+,-,+,\cdots$.}
% segment-id: o014.li2.chapter3.complexes-over-additive-categories.q002
	\[ \left(X[n]\right)^k:=X^{k+n},\quad
	d_{X[n]}^k:=(-1)^nd_X^{k+n}. \]
	This construction gives an additive functor
	$[n]:\cate{C}(\mathcal{A})\to\cate{C}(\mathcal{A})$. If $f:X\to Y$ is a
	morphism of complexes, then $f[n]:X[n]\to Y[n]$ is given by
	$f[n]^k:=f^{n+k}$.
\end{definition}

% segment-id: o014.li2.chapter3.complexes-over-additive-categories.p007
It is clear that $X[n]$ is again a complex, so the functor is well defined.
The following properties are also immediate.
% segment-id: o014.li2.chapter3.complexes-over-additive-categories.l002
\begin{itemize}
% segment-id: o014.li2.chapter3.complexes-over-additive-categories.i004
	\item $[0]=\identity_{\cate{C}(\mathcal{A})}$.
% segment-id: o014.li2.chapter3.complexes-over-additive-categories.i005
	\item $[n][m]=[n+m]$. Thus $[n]$ is an automorphism of
	$\cate{C}(\mathcal{A})$, with inverse $[-n]$.
% segment-id: o014.li2.chapter3.complexes-over-additive-categories.i006
	\item The family $(d_X^n)_{n\in\Z}$ gives a morphism of complexes
	$d_X:X\to X[1]$ (Hint: $d_X^{n+1}d_X^n=0$).
\end{itemize}

% segment-id: o014.li2.chapter3.complexes-over-additive-categories.c001
\begin{convention}\label{con:concentrated-cplx}
	Every object $S$ of $\mathcal{A}$ may be regarded as a complex
	$(S^n,d_S^n)_{n\in\Z}$ with $S^0:=S$, every other term equal to $0$, and
	every $d_S^n$ the zero morphism. This identification gives a fully faithful
	additive functor $\mathcal{A}\hookrightarrow\cate{C}(\mathcal{A})$. More
	generally, for every $n\in\Z$, $S[-n]$ is the complex concentrated in
	degree $n$: it places $S$ in degree $n$ and $0$ in every other degree.
\end{convention}

% segment-id: o014.li2.chapter3.complexes-over-additive-categories.p008
The following result is immediate.

% segment-id: o014.li2.chapter3.complexes-over-additive-categories.pr002
\begin{proposition}\label{prop:CA-functoriality}
	\index[sym1]{CF@$\cate{C}F$}
	For an additive functor $F:\mathcal{A}\to\mathcal{A}'$, the assignment
	$(X^n,d_X^n)_n\mapsto(FX^n,Fd_X^n)_n$ defines a functor
	$\cate{C}F:\cate{C}(\mathcal{A})\to\cate{C}(\mathcal{A}')$. It commutes
	with translation functors:
	$[m]\circ\cate{C}F=\cate{C}F\circ[m]$ for every $m\in\Z$.
\end{proposition}

% segment-id: o014.li2.chapter3.complexes-over-additive-categories.r001
\begin{remark}\label{rem:chain-cplx-translation}
	\index{translation functor}
	\index[sym1]{[n]}
	Every result in this section also applies to the chain complexes
	$(X_\bullet,d_\bullet)$ mentioned in Remark~\ref{rem:cochain-vs-chain}.
	The translation functor on the category of chain complexes is defined by
% segment-id: o014.li2.chapter3.complexes-over-additive-categories.q003
	\[ X[n]_k:=X_{k+n},\quad d^{X[n]}_k=(-1)^nd^X_{n+k} \]
	If one passes between the two conventions via $X_n=X^{-n}$ and
	$d_n=d^{-n}$, then the translation $[1]$ for chain complexes corresponds
	to the translation $[-1]$ for cochain complexes.
\end{remark}
